\subsection{Levantamento de Requisitos}

O levantamento de requisitos teve como objetivo definir as funcionalidades 
e as propriedades de qualidade necessárias à integração do módulo de 
solicitações ao sistema Banco de Especialistas. A solução foi concebida 
para preservar o comportamento das funcionalidades existentes e formalizar 
o processo de alteração de categoria e de dados cadastrais dos especialistas, 
mediante validação das entradas, controle de aprovação e registro das 
operações realizadas.

Os requisitos funcionais especificam as operações disponibilizadas pelo 
módulo e seus efeitos sobre os registros do sistema. Os requisitos não 
funcionais estabelecem as condições de integridade, segurança, 
rastreabilidade, compatibilidade e manutenção que orientam a implementação 
da solução.

\subsubsection{Requisitos Funcionais}

Os requisitos funcionais abrangem a criação e o processamento de 
solicitações, a consulta e o acompanhamento das demandas, a aplicação 
das alterações aprovadas, a preservação do estado anterior dos registros 
e o controle de acesso às operações. As decisões e as alterações devem 
ser executadas exclusivamente por usuários autorizados, conforme as 
regras de negócio estabelecidas.

\FloatBarrier
\begin{quadro}[htbp]
\centering
\caption[Requisitos funcionais do módulo]{Requisitos funcionais do módulo de solicitações.}
\label{qua:rf}
\small
\setlength{\tabcolsep}{4pt}
\begin{tabular}{p{1cm}p{6.5cm}p{2.2cm}p{2.3cm}p{2cm}}
\toprule
\textbf{ID} & \textbf{Descrição} & \textbf{Tipo de solicitação} & \textbf{Origem} & \textbf{Implementação} \\
\midrule
RF01 & Permitir a criação de solicitações de mudança de categoria, com indicação da categoria de destino e justificativa obrigatória. & Categoria & Seção~\ref{sec:fluxo-antigo} & Seção~\ref{sec:modulo-solicitacoes} \\
RF02 & Permitir a criação de solicitações de atualização de pontuação, com indicação dos itens e das respectivas comprovações. & Pontuação & Seção~\ref{sec:fluxo-antigo} & Seção~\ref{sec:modulo-solicitacoes} \\
RF03 & Permitir a criação de solicitações de complemento cadastral ou de recurso, com descrição do problema identificado. & Complemento & Seção~\ref{sec:fluxo-antigo} & Seção~\ref{sec:modulo-solicitacoes} \\
RF04 & Permitir a consulta das solicitações, com filtros por tipo e por situação, conforme o perfil do usuário (solicitante ou avaliador). & Todos & Seção~\ref{sec:fluxo-antigo} & Seção~\ref{sec:interfaces-web} \\
RF05 & Permitir a aprovação ou rejeição de solicitações, registrando o responsável e a decisão tomada. & Todos & Seção~\ref{sec:fluxo-antigo} & Seção~\ref{sec:modulo-solicitacoes} \\
RF06 & Aplicar, após a aprovação, as alterações correspondentes à categoria de destino do especialista. & Categoria & Seção~\ref{sec:regras-negocio} & Seção~\ref{sec:modulo-solicitacoes} \\
RF07 & Aplicar, após a aprovação, a atualização da pontuação do especialista correspondente aos itens aprovados. & Pontuação & Seção~\ref{sec:regras-negocio} & Seção~\ref{sec:modulo-solicitacoes} \\
RF08 & Registrar o estado anterior do especialista antes da aplicação de alterações registradas em snapshot. & Categoria & Seção~\ref{sec:fluxo-antigo} & Seção~\ref{sec:snapshot-auditoria} \\
RF09 & Permitir a reversão de alterações aprovadas, restaurando o estado anterior registrado. & Categoria & Seção~\ref{sec:fluxo-antigo} & Seção~\ref{sec:snapshot-auditoria} \\
RF10 & Registrar e disponibilizar uma linha do tempo cronológica dos eventos associados a cada solicitação. & Todos & Seção~\ref{sec:regras-negocio} & Seção~\ref{sec:snapshot-auditoria} \\
RF11 & Restringir as operações de criação, consulta e decisão a usuários autenticados e autorizados, conforme o perfil de acesso. & Todos & Seção~\ref{sec:fluxo-antigo} & Seção~\ref{sec:modulo-solicitacoes} \\
\bottomrule
\end{tabular}
\fonte{Elaboração própria (2026).}
{\footnotesize\noindent\textbf{Nota:} O requisito RF03 refere-se à abertura 
de solicitações de complemento cadastral ou recurso. Após a aprovação ou 
rejeição da demanda, o sistema registra a decisão e o encaminhamento 
correspondente, sem alterar automaticamente o perfil do especialista.}
\end{quadro}
\FloatBarrier

\subsubsection{Requisitos Não Funcionais}

Os requisitos não funcionais estabelecem as propriedades de qualidade 
necessárias ao funcionamento do módulo e à sua integração ao sistema 
existente. Esses requisitos abrangem a integridade das operações, a 
rastreabilidade das alterações, a segurança do controle de acesso, a 
validação dos dados, a compatibilidade com as funcionalidades existentes, 
a extensibilidade da solução, a testabilidade das regras de negócio e a 
reprodutibilidade do ambiente de testes.

\FloatBarrier
\begin{quadro}[htbp]
\centering
\caption[Requisitos não funcionais do módulo]{Requisitos não funcionais do módulo de solicitações.}
\label{qua:rnf}
\small
\setlength{\tabcolsep}{4pt}
\begin{tabular}{p{1cm}p{2.3cm}p{6.5cm}p{3cm}}
\toprule
\textbf{ID} & \textbf{Categoria} & \textbf{Descrição} & \textbf{Mecanismo de atendimento} \\
\midrule
RNF01 & Integridade & As alterações decorrentes da aprovação devem ser executadas atomicamente, evitando a persistência parcial dos registros envolvidos na operação. & Transações PostgreSQL e SQLAlchemy \\
RNF02 & Rastreabilidade & As informações de auditoria devem ser preservadas de forma que seja possível reconstituir, a qualquer momento, o histórico de decisões e o estado anterior das alterações registradas em snapshot. & Entidade Snapshot e linha do tempo de eventos \\
RNF03 & Segurança & O controle de acesso deve ser aplicado no servidor, de modo que as operações não possam ser executadas por chamadas diretas à API que contornem a interface. & Verificação de permissões nas rotas da API \\
RNF04 & Validação & Os dados recebidos pela API devem ser validados no servidor, de modo que a validação do servidor seja autoritativa, independentemente das verificações realizadas pela interface. & Schemas Marshmallow \\
RNF05 & Compatibilidade & A incorporação do módulo deve preservar o comportamento das funcionalidades existentes não abrangidas pelo módulo. & Blueprints Flask e estratégia de mudança aditiva (\textit{Additive Change}) \\
RNF06 & Extensibilidade & A inclusão de novos tipos de solicitação deve permitir o reaproveitamento do fluxo principal, limitando as alterações necessárias às regras específicas de cada tipo. & Formulário único e configuração \texttt{SCHEMA\_POR\_TIPO} \\
RNF07 & Testabilidade & As regras de negócio devem permitir testes isolados sem requisição HTTP, de modo que possam ser verificadas independentemente da camada de comunicação. & Separação da camada de serviço e Pytest \\
RNF08 & Reprodutibilidade & O ambiente de testes deve poder ser configurado e executado de maneira reproduzível, com dependências e procedimentos documentados. & Docker e Docker Compose \\
\bottomrule
\end{tabular}
\fonte{Elaboração própria (2026).}
\end{quadro}
\FloatBarrier